\Needspace{9\baselineskip}\begin{center}\itshape Weitere Untersuchungen über die Perioden der Wurzeln.\\ Summen mit einer geraden Anzahl von Gliedern sind reelle Größen.\end{center}


\Needspace{6\baselineskip}\Article{355}
Da \(n\) stets als ungerade vorausgesetzt wird, gehört \(2\) zu den Faktoren von \(n-1\), und die Menge \(\Omega\) besteht aus \(\frac12(n-1)\) Perioden mit zwei Gliedern. Eine solche Periode wie \((2,\lambda)\) besteht aus den Wurzeln \([\lambda]\) und \([\lambda g^{\frac12(n-1)}]\), wobei \(g\) wie oben eine beliebige primitive Wurzel für den Modul \(n\) bezeichnet. Es gilt aber \(g^{\frac12(n-1)}\equiv-1\ (\mathrm{mod.}\,n)\) und daher \(\lambda g^{\frac12(n-1)}\equiv-\lambda\) (siehe Artikel 62), woraus \([\lambda g^{\frac12(n-1)}]=[-\lambda]\) folgt. Setzt man also \([\lambda]=\cos\frac{kP}{n}+i\sin\frac{kP}{n}\) und daher \([-\lambda]=\cos\frac{kP}{n}-i\sin\frac{kP}{n}\), so ergibt sich die Summe \((2,\lambda)=2\cos\frac{kP}{n}\). Daraus ziehen wir hier nur den Schluss, dass der Wert jeder Summe mit zwei Gliedern eine reelle Größe ist. Da sich jede Periode, deren Gliederzahl gerade \(=2a\) ist, in \(a\) Perioden mit je zwei Gliedern zerlegen lässt, ist allgemeiner klar, dass der Wert jeder Summe,
\SourceLeaf{446}
deren Gliederzahl gerade ist, stets eine reelle Größe ist. Wird also in Artikel 352 unter den Faktoren \(\alpha,\allowbreak \beta,\allowbreak \gamma\) usw. der Faktor zwei für die letzte Stelle aufgespart, so werden sämtliche Operationen bis zu den Summen mit zwei Gliedern durch reelle Größen ausgeführt; imaginäre Größen werden erst dann eingeführt, wenn man von diesen Summen zu den Wurzeln selbst übergeht.

\Needspace{9\baselineskip}\begin{center}\itshape Über die Gleichung, die die Verteilung der Wurzeln \(\Omega\) in zwei Perioden bestimmt.\end{center}


\Needspace{6\baselineskip}\Article{356}
Die Hilfsgleichungen, durch die für jeden Wert von \(n\) die Summen bestimmt werden, welche die Menge \(\Omega\) bilden, verdienen höchste Aufmerksamkeit: Sie hängen auf bemerkenswerte Weise mit den verborgensten Eigenschaften der Zahl \(n\) zusammen. Hier wollen wir die Untersuchung jedoch auf die folgenden zwei Fälle beschränken: \textit{Erstens} behandeln wir die quadratische Gleichung, deren Wurzeln die Summen mit \(\frac12(n-1)\) Gliedern sind; \textit{zweitens}, für den Fall, dass \(n-1\) den Faktor \(3\) enthält, die kubische Gleichung, deren Wurzeln die Summen mit \(\frac13(n-1)\) Gliedern sind.

Schreiben wir der Kürze wegen \(m\) für \(\frac12(n-1)\) und bezeichnen wir mit \(g\) eine beliebige primitive Wurzel für den Modul \(n\), so besteht die Menge \(\Omega\) aus zwei Perioden \((m,1)\) und \((m,g)\); die erste enthält die Wurzeln \([1],\allowbreak [gg],\allowbreak [g^4]\ldots[g^{n-3}]\), die zweite \([g],\allowbreak [g^3],\allowbreak [g^5]\ldots[g^{n-2}]\). Die kleinsten positiven Reste der Zahlen \(gg,\allowbreak g^4\ldots g^{n-3}\) nach dem Modul \(n\) seien in beliebiger Reihenfolge \(R,\allowbreak R',\allowbreak R''\) usw.; ebenso seien die Reste von \(g,\allowbreak g^3,\allowbreak g^5\ldots g^{n-2}\) gleich \(N,\allowbreak N',\allowbreak N''\) usw. Dann stimmen die Wurzeln, aus denen \((m,1)\) besteht, mit \([1],\allowbreak [R],\allowbreak [R'],\allowbreak [R'']\) usw. überein und die Wurzeln der Periode \((m,g)\) mit \([N],\allowbreak [N'],\allowbreak [N'']\) usw. Nun ist klar, dass sämtliche Zahlen \(1,\allowbreak R,\allowbreak R',\allowbreak R''\) usw. \textit{quadratische Reste} der Zahl \(n\) sind; da sie sämtlich verschieden und kleiner als \(n\) sind und ihre Anzahl \(=\frac12(n-1)\), also gleich der Anzahl aller positiven Reste von \(n\) unterhalb \(n\) ist, stimmen diese Reste genau mit jenen Zahlen überein. Daraus folgt unmittelbar, dass sämtliche Zahlen \(N,\allowbreak N',\allowbreak N''\) usw., die sowohl untereinander als auch von \(1,\allowbreak R,\allowbreak R'\) usw. verschieden sind und zusammen mit diesen alle Zahlen \(1,\allowbreak 2,\allowbreak 3\ldots n-1\) erschöpfen, mit sämtlichen positiven \textit{quadratischen Nichtresten} von \(n\) unterhalb \(n\) übereinstimmen müssen. Setzt man nun voraus, dass die Gleichung, deren Wurzeln die Summen \((m,1),\allowbreak (m,g)\) sind, \[xx-Ax+B=0\] lautet, so gilt \[A=(m,1)+(m,g)=-1,\quad B=(m,1)\times(m,g)\]
\SourceLeaf{447}

Das Produkt aus \((m,1)\) und \((m,g)\) ist nach Artikel 345 \[=(m,N+1)+(m,N'+1)+(m,N''+1)+\text{usw.}=W\] und lässt sich daher auf eine Form wie \(\alpha(m,0)+\beta(m,1)+\gamma(m,g)\) bringen. Zur Bestimmung der Koeffizienten \(\alpha,\allowbreak \beta,\allowbreak \gamma\) bemerken wir \textit{erstens}, dass \(\alpha+\beta+\gamma=m\) gilt (da die Anzahl der Summen in \(W\) \(=m\) ist); \textit{zweitens}, dass \(\beta=\gamma\) gilt (dies folgt aus Artikel 350, weil das Produkt \((m,1)\times(m,g)\) eine symmetrische Funktion der Summen \((m,1),\allowbreak (m,g)\) ist, aus denen die größere Summe \((n-1,1)\) besteht); \textit{drittens}: Da sämtliche Zahlen \(N+1,\allowbreak N'+1,\allowbreak N''+1\) usw. strikt zwischen den Grenzen \(2\) und \(n+1\) liegen, ist klar, dass entweder keine Summe in \(W\) auf \((m,0)\) zurückgeführt wird und daher \(\alpha=0\) gilt, wenn unter den Zahlen \(N,\allowbreak N',\allowbreak N''\) usw. die Zahl \(n-1\) nicht vorkommt; oder genau eine, nämlich \((m,n)\), sodass \(\alpha=1\) gilt, wenn \(n-1\) unter den Zahlen \(N,\allowbreak N',\allowbreak N''\) usw. vorkommt. Daraus schließen wir, dass im ersten Fall \(\alpha=0,\allowbreak \ \beta=\gamma=\frac12m\), im zweiten \(\alpha=1,\allowbreak \ \beta=\gamma=\frac12(m-1)\) gilt. Zugleich folgt daraus, da die Zahlen \(\beta\) und \(\gamma\) notwendig ganz sind, dass der erste Fall eintritt, also \(n-1\) (oder, was dasselbe ist, \(-1\)) unter den Nichtresten von \(n\) nicht vorkommt, wenn \(m\) gerade, also \(n\) von der Form \(4k+1\) ist; der zweite Fall tritt dagegen ein, also \(n-1\) oder \(-1\) kommt unter den Nichtresten von \(n\) vor, wenn \(m\) ungerade, also \(n\) von der Form \(4k+3\) ist\HistoricalNote{*)}{Auf diese Weise haben wir einen neuen Beweis des Satzes erhalten, dass \(-1\) ein Rest aller Primzahlen von der Form \(4k+1\) und ein Nichtrest aller der Form \(4k+3\) ist, der oben (Artikel 108, 109 und 262) bereits auf mehrere verschiedene Arten bewiesen wurde. Will man diesen Satz lieber voraussetzen, so braucht man zur Unterscheidung der beiden verschiedenen Fälle die Bedingung nicht zu berücksichtigen, dass \(\beta,\allowbreak \gamma\) schon an sich ganze Zahlen sind.}. Daher ist das gesuchte Produkt wegen \((m,0)=m\), \((m,1)+(m,g)=-1\) im ersten Fall \(=-\frac12m\), im zweiten \(=\frac12(m+1)\), und die gesuchte Gleichung lautet in jenem Fall \(xx+x-\frac14(n-1)=0\), deren Wurzeln \(-\frac12\pm\frac12\surd n\) sind, in diesem dagegen \(xx+x+\frac14(n+1)=0\), deren Wurzeln \(-\frac12\pm\frac12i\surd n\) sind.

Welche Wurzel aus \(\Omega\) also auch für \([1]\) gewählt wurde, die Differenz zwischen den Summen \(\Sigma[\mathfrak R]\) und \(\Sigma[\mathfrak N]\), wobei für \(\mathfrak R\) sämtliche positiven quadratischen Reste und für \(\mathfrak N\) sämtliche positiven quadratischen Nichtreste von \(n\) unterhalb \(n\) einzusetzen sind, ist \(=\pm\surd n\) für \(n\equiv1\) und \(=\pm i\surd n\) für \(n\equiv3\ (\mathrm{mod.}\,4)\). Ebenso folgt daraus leicht, wenn \(k\) eine beliebige ganze, durch \(n\) nicht teilbare Zahl bezeichnet, dass {\postdisplaypenalty=10000 \[\begin{gathered}\Sigma\cos\frac{k\mathfrak RP}{n}-\Sigma\cos\frac{k\mathfrak NP}{n}=\pm\surd n\quad\text{und}\\\Sigma\sin\frac{k\mathfrak RP}{n}-\Sigma\sin\frac{k\mathfrak NP}{n}=0\end{gathered}\]} für \(n\equiv1\ (\mathrm{mod.}\,4)\) gilt; dagegen ist für \(n\equiv3\ (\mathrm{mod.}\,4)\) jene Differenz \(=0\), diese
\SourceLeaf{448}
\(=\pm\surd n\). Diese Sätze sind wegen ihrer Eleganz sehr bemerkenswert. Übrigens bemerken wir, dass die oberen Vorzeichen immer gelten, wenn für \(k\) die Einheit oder allgemeiner ein quadratischer Rest von \(n\) genommen wird, die unteren, wenn für \(k\) ein Nichtrest gewählt wird, und dass sich diese Sätze unter Wahrung oder vielmehr Steigerung ihrer Eleganz auch auf beliebige zusammengesetzte Werte von \(n\) ausdehnen lassen. Über diese Dinge, die einer tieferen Untersuchung bedürfen, müssen wir hier jedoch schweigen und ihre Betrachtung einer anderen Gelegenheit vorbehalten

